% introduction.tex -- Introduction
% Three-paragraph structure: (1) problem, (2) landscape, (3) contribution.

Attention is the central computational primitive of modern deep learning~\cite{vaswaniAttentionAllYou2017}. Given a query, it computes a softmax-weighted average of stored values, a powerful operation but a fundamentally \emph{deterministic} one. The same query always produces the same output. Attention retrieves; it does not generate. Yet generation from a structured memory is precisely what many downstream tasks require: producing novel but plausible continuations, interpolating between stored prototypes, or exploring the space of patterns consistent with partial evidence. A natural question arises: \emph{can the attention mechanism itself be made stochastic, in a principled way, so that it samples from the space of memories rather than merely returning their weighted average?}

The ingredients for an answer already exist, though they have not been assembled. On one side, a line of work on associative memory, from Hopfield networks~\cite{hopfieldNeuralNetworksPhysical1982} through dense associative memories~\cite{krotovDenseAssociativeMemory2016,krotovLargeAssociativeMemory2021} to the modern continuous formulation of Ramsauer et al.~\cite{ramsauerHopfieldNetworksAll2021}, has revealed that each attention head implicitly performs gradient descent on a smooth, confining energy whose minima are stored patterns. On the other side, a mature theory of Langevin dynamics~\cite{robertsTweedieExponentialConvergence1996,wellingBayesianLearningStochastic2011,durmusNonasymptoticAnalysisUnadjusted2017} shows how to convert \emph{any} such energy into a sampler for the corresponding Boltzmann distribution: one simply adds calibrated noise to the gradient update. Meanwhile, energy-based models~\cite{lecunTutorialEnergyBased2006,songHowTrainEnergyBased2021} and score-based diffusion methods~\cite{songGenerativeModelingEstimating2019,songScoreBasedGenerativeModeling2021} have demonstrated the power of Langevin-type sampling for generation, but rely on black-box neural networks whose scores must be learned. The Energy Transformer~\cite{hooverEnergyTransformer2024} brings Hopfield energies into a deep architecture, yet remains a discriminative model that descends the energy rather than sampling from it. The classical retrieval-generation duality (Hopfield networks retrieve, Boltzmann machines sample, and both are defined by the same quadratic interaction energy over binary states~\cite{ackeleyLearningAlgorithmBoltzmann1985}) has not been lifted to the modern continuous setting.

In this paper, we close that gap. We show that applying the unadjusted Langevin algorithm to the modern Hopfield energy yields a \emph{stochastic attention} update: a single iteration comprising a contraction toward the origin, a softmax attention pull toward stored memories, and an isotropic Gaussian perturbation whose magnitude is governed by the temperature. The inverse temperature $\beta$ interpolates continuously between exact retrieval ($\beta\to\infty$) and open-ended exploration ($\beta\to 0$), providing a principled generation mechanism that requires no learned score network, no training loop, and no contrastive objective. Because the gradient of the Hopfield energy is exactly the identity minus the attention map, every step of the sampler is computed by the same primitive operations as a standard attention head. In particular, $\mathbf{X}$ can be the key matrix of any pretrained attention layer, making stochastic attention a zero-shot stochastic decoding layer compatible with retrieval-augmented generation and in-context learning settings. The energy's analytic structure (infinite differentiability, Lipschitz gradients, and a quadratic confining bound) delivers convergence guarantees that generic energy-based models cannot offer without additional assumptions. In practice, the same algorithm realizes two qualitatively distinct operating regimes: at high $\beta$ it performs \emph{structured retrieval}, generating outputs that closely resemble stored patterns; once $\beta$ falls below a signal-to-noise threshold ($\mathrm{SNR}\approx0.025$, derived in Section~4), it performs \emph{genuine generation}, producing novel outputs that outperform learned baselines on both novelty and diversity. Section~2 surveys related work; Section~3 reviews the necessary background on attention, modern Hopfield networks, and Langevin dynamics; Section~4 derives the stochastic attention update, presents the algorithm, and characterizes its properties; Section~5 validates both regimes experimentally across four domains.
